% Vista pública derivada de 73_rev2_electron_lectores_torres.tex: cuerpo exacto y único retítulo académico del lector temporal de 108 estados. El fragmento del handoff permanece byte-exacto.
\section{Complete construction of the beta electron: central route, signature, character, and governing mass}
\Needspace{7\baselineskip}
\subsection{Central route of the electron and sequence of genealogical records}
The electron occupies the unique center \((5,5)\), with active orientation \(++\):

\[
 \Gamma_e=(5,5,++),\qquad
 \tau_{12}=+++---+++---.
\]
Their displacement readers are

\[
 (D_1,D_3,D_4,D_{27},D_{36})=(54,56,68,48,32).
\]
They must not be merged:

\begin{enumerate}
\item the raw signature \((24,0,12,0,-12)\);
\item the even CT--108 signature \((-12,-4,12,-6,12)\);
\item the vector of twelve events;
\item the affine origin \(v_e=0\) of the relative character.
\end{enumerate}
Nor should \(\Gamma_e\) be confused with the word \(w_e=201101\) of the reader of Euler's constant \(e\).

\Needspace{7\baselineskip}
\subsection{Basal electron selector and transition to the beta ruling value}
The internal selector, constructed without using the observed mass, closes

\[
 -22=-\left|((\mathbb Z/9\mathbb Z)\times\mathbb F_3)\setminus\iota(R_\pi)\right|,
 \qquad
 +15=|R_\pi\times\mathbb F_3|.
\]
The reading of vacancies also contributes

\[
 V_e=\begin{pmatrix}21&22\\6&7\end{pmatrix},\qquad
 \det V_e=15.
\]
The sign negativo of 22 records deficit/cut; the positive of 15, retencion/arean oriented. The matrix is local to the electron and does not  extiende as selector universal.

The basal stage is

\[
 e_0=\frac{\sqrt3}{4}
 \left[\exp\!\left(\frac{\varphi}{\pi^2}\right)
 -22\alpha_{\rm HMT}^3\right](1+15\Delta_4),
\]
\[
 e_0=0.510998922488\ldots\ \mathrm{MeV}.
\]
The entire genealogies are

\[
 22=\frac{396}{18}=27-5=24-2,
 \qquad
 15=\frac{270}{18}=\binom62=\binom64.
\]
\Needspace{7\baselineskip}
\subsection{Bidirectional readers and beta}
For every route,

\[
 K_D(\Gamma)=\left(D_{27}+D_{36},D_1,\frac{D_4-D_3}{2}\right),
\]
\[
 K_\Omega(\Gamma)=(\Omega_{90,120},54,P_6).
\]
Across the 324 routes there appear 14 profiles \(K_D\), 9 profiles \(K_\Omega\), 40 joint pairs, and 18 coincidences. The only common value is

\[
 (80,54,6).
\]
This 80 comes from the bidirectional reader of the electronic route, not from period 80 of the Euler word. Therefore

\[
 \kappa_e=80-54\alpha_{\rm HMT}+6\Delta_4,
\]
\[
 \boxed{
 \mathfrak e_e^\beta=e_0R_{\rm act}^{\kappa_e}
 =0.5109989506900008086082571648\ldots\ \mathrm{MeV}.}
\]
This is the governing electronic output. The basal output does not replace it. Equality with the antiparticle mass proceeds from route conjugation and from the parity of the mass operator, not from adding a second numerical row.

\Needspace{7\baselineskip}
\section{Physical realization by spectral coupling}
\Needspace{7\baselineskip}
\subsection{Obstruction to the equivariant punctual selection of a non-central multisection}
Let \(F\subset\mathcal R_{324}\) be a route fiber compatible with a structural class. Its linear space is

\[
 \mathcal H_F=\bigoplus_{\Gamma\in F}\mathbb C e_\Gamma.
\]
The central character and the readers define diagonal operators:

\[
 \widehat{\mathcal A}_F^{(0)}e_\Gamma
 =\chi_{\rm full}(\sigma_\Gamma,\eta_\Gamma)e_\Gamma,
\]
\[
 \widehat K_F^r e_\Gamma=\kappa_r(\Gamma)e_\Gamma,
 \qquad r\in\{D,\Omega\}.
\]
The two--way correction is

\[
 \widehat{\mathcal M}_F^{\beta,r}
 =R_{\rm act}^{\widehat K_F^r/2}
  \widehat{\mathcal M}_F^{(0)}
  R_{\rm act}^{\widehat K_F^r/2}.
\]
When \(\dim\mathcal H_F>1\), these operators generally have several eigenvalues. Choosing one because it approximates a known mass would reverse causality. Nor is it correct to require a physical name to identify a unique cell: the noncentral species resides in a representation of the fiber.

\Needspace{7\baselineskip}
\subsection{Spectral functor of the physical realization: positive coupling kernel, persistent subspace, and spectral measure}
A physical preparation \(P\) provides representation data, not a mass:

\[
 \mathfrak r(P)=
 (q,\,j,\,\epsilon_{2\pi},\,\epsilon_{4\pi},\,
  \text{statistics},\,\text{generation},\,\text{flavor},\,
  \text{color orbit},\,\text{interaction channel}).
\]
Let \(\mathfrak M\) be a coupling compatible with that representation. A positive kernel is defined

\[
 K_{P,\mathfrak M}:F\times F\longrightarrow\mathbb C
\]
which satisfies:

\begin{enumerate}
\item \(K=K^*\) y \(K\ge0\);
\item \(K(\Gamma,\Gamma')=0\) if the routes violate charge, parity, colour, statistics, or the
\(P\) channel;

\item \(K\) commutes with the symmetries that the coupling does not break;
\item \(K(C_M\Gamma,C_M\Gamma')=\overline{K(\Gamma,\Gamma')}\) under conjugation;
\item its entries are computed from the genealogical ledger, boundary, memory, carry, and holonomy, without
external magnitudes.

\end{enumerate}
The coupled evolution operator is

\[
 \mathcal U_{P,\mathfrak M}
 =K_{P,\mathfrak M}^{1/2}\,\mathcal U_{\rm TPK}\,
  K_{P,\mathfrak M}^{1/2}.
\]
The persistent subspace is obtained through the spectral projection onto the modes that survive return and pruning:

\[
 \Pi_{P,\mathfrak M}^{\rm surv}
 =\mathbf1_{\mathcal S_{\rm surv}}(\mathcal U_{P,\mathfrak M}).
\]
The correct physical realization is the spectral functor.

\[
 \boxed{
 \mathcal S_{\rm phys}^{\rm spec}(P,\mathfrak M)
 =\left(
 \mathcal H_{P,\mathfrak M}^{\rm surv},
 \widehat{\mathcal M}_{P,\mathfrak M},
 \mu_{P,\mathfrak M}
 \right),}
\]
where

\[
 \mathcal H_{P,\mathfrak M}^{\rm surv}
 =\Pi_{P,\mathfrak M}^{\rm surv}\mathcal H_F,
\]
\[
 \widehat{\mathcal M}_{P,\mathfrak M}
 =\Pi_{P,\mathfrak M}^{\rm surv}
  \widehat{\mathcal M}_F
  \Pi_{P,\mathfrak M}^{\rm surv},
\]
and \(\mu_{P,\mathfrak M}\) is its spectral measure in the prepared state.

When the coupling is specified by an irreducible representation \(\lambda\) of a finite group \(G_P\), the corresponding projector is obtained without choosing a distinguished route:

\[
 P_{P,\lambda}
 =\frac{\dim\lambda}{|G_P|}
  \sum_{g\in G_P}\overline{\chi_\lambda(g)}\,U(g).
\]
The bidirectional output is then the compressed pair.

\[
 \mathbb M(P,\lambda)=
 \left(
 P_{P,\lambda}\widehat{\mathcal M}^{\beta,D}P_{P,\lambda},
 P_{P,\lambda}\widehat{\mathcal M}^{\beta,\Omega}P_{P,\lambda}
 \right).
\]
If both operators commute with the irreducible representation, Schur's lemma forces their scalarity on the block. If they do not commute, the correct result is the spectrum of the minimal invariant block.

\Needspace{7\baselineskip}
\subsection{Naturality of the Spectral Functor under Intertwiners of the TPK Dynamics, Coupling, and Mass Operator}
If \(V:F\to F'\) intertwines the TPK dynamics, the coupling kernel, and the mass operators, then

\[
 V\Pi_{P,\mathfrak M}^{\rm surv}
 =\Pi_{P',\mathfrak M'}^{\rm surv}V,
 \qquad
 V\widehat{\mathcal M}_{P,\mathfrak M}
 =\widehat{\mathcal M}_{P',\mathfrak M'}V.
\]
By functional calculus, the spectral measure is transported by \(V\). The output does not depend on a numbering of the cells or on the choice of a basis within an orbit. This is the condition that must replace every non-equivariant nominal selector.

\Needspace{7\baselineskip}
\subsection{Conditions for the existence of a scalar mass under the dimensional reader}
A scalar mass is determined in any of the following cases:

\begin{enumerate}
\item the persistent subspace has rank one;
\item an irreducible representation fixes a unique spectral projector;
\item the measurement protocol defines a normalized positive state
\(\omega_{P,\mathfrak M}\) on the algebra of observables;

\item a resonance defines an isolated pole of the resolvent.
\end{enumerate}
In the third case,

\[
 m(P,\mathfrak M)
 =\omega_{P,\mathfrak M}
  (\widehat{\mathcal M}_{P,\mathfrak M}).
\]
In the absence of one of these mechanisms, the scientific output is the operator, its spectrum, and its multiplicities. That output is not incomplete: it is the object determined by the dynamics.

\Needspace{7\baselineskip}
\subsection{Finite construction of the mass kernel by means of representation, persistence, and coupling filters, followed by irreducible normalization}
The kernel is obtained in four finite stages:

\begin{enumerate}
\item \textbf{Representation filter.} Retain the routes whose sector, charge branch,
color, holonomy \(2\pi/4\pi\), parity, and statistics match \(P\).

\item \textbf{Persistence filter.} The partial order formed by class return,
orbit return, cell return, prefix separation, carry debt, and decimal stability is used. The result may be a Pareto orbit; uniqueness is not imposed.

\item \textbf{Coupling filter.} Boundary incidence and the
interaction channel are evaluated. Decoupled routes receive zero weight.

\item \textbf{Normalization.} On each irreducible component the trace of the
kernel is normalized to one.

\end{enumerate}
The oriented atlas provides all data of the first two stages. The third requires each physical realization to declare its channel; the fourth is canonical once the irreducible components have been fixed.

\Needspace{7\baselineskip}
\section{Genealogical Transducer of Towers}
\Needspace{7\baselineskip}
\subsection{Pair of associated words for the same genealogical route}
Each \(\Gamma\in\mathcal R_{324}\) determines simultaneously a word temporal

\[
 w_\Gamma\in\mathbb F_3^{108}
\]
and a dodecaphonic word signed

\[
 \tau_\Gamma\in\{-1,0,+1\}^{12}.
\]
The first preserves displacements along CT--108; the second preserves the accumulated orientation of the twelve events. The two words are linked by the same genealogical record, but they are not interchangeable. The refinement of towers reads both and also preserves the sheet, carry, winding, and return memory.

\Needspace{7\baselineskip}
\subsection{Correspondence between harmonic height and temporal length}
Write

\[
 \mathfrak S=30\langle3,4\rangle
 =\{90,120\}\cup\{30r:r\ge6\}.
\]
For \(h=30r\), the associated temporal length is

\[
 \ell_h=9r.
\]
In particular,

\[
 90\longleftrightarrow27,\qquad
 120\longleftrightarrow36,\qquad
 360\longleftrightarrow108.
\]
This correspondence does not identify the tower with time: it intertwines the height of the harmonic operator with the length of the route segment that generates it.

\Needspace{7\baselineskip}
\subsection{Temporal cycle shift reader of 108 states}
We define the cyclic discrepancy

\[
 D_\Gamma(n)=
 \sum_{t=0}^{107}
 \mathbf1_{\{w_{\Gamma,t}\ne
 w_{\Gamma,t+n\!\!\!\pmod{108}}\}}.
\]
The relative coefficient of the temporal branch is

\[
 \boxed{\eta_{30r}^{D}(\Gamma)
 =D_\Gamma(9r)-D_{\Gamma_e}(9r).}
\]
Subtraction does not use a mass as an origin: it uses the central route as the affine origin of the relative character. It holds.

\[
 \eta_{90}^{D}+\eta_{120}^{D}
 =[D_\Gamma(27)+D_\Gamma(36)]-80,
\]
so that the first layer exactly prolongs the first component of the reader \(K_D\).

\Needspace{7\baselineskip}
\subsection{Dodecaphase Accumulation Reader}
For \(r\ge1\), let

\[
 A_\Gamma(r)=
 \sum_{j=0}^{11}
 \left(
 \sum_{u=0}^{r-1}
 \tau_{\Gamma,j+u\!\!\!\pmod{12}}
 \right)^2.
\]
The relative coefficient of the oriented branch is

\[
 \boxed{\eta_{30r}^{\Omega}(\Gamma)
 =A_\Gamma(r)-A_{\Gamma_e}(r).}
\]
Then

\[
 4\eta_{90}^{\Omega}-3\eta_{120}^{\Omega}
 =\Omega(\tau_\Gamma)-80,
\]
so this layer extends the first component of \(K_\Omega\). The two families form the bidirectional lift

\[
 \mathcal F_{\rm tower}^{2w}:
 \mathcal R_{324}\longrightarrow
 \mathcal E_{\mathfrak S}\times\mathcal E_{\mathfrak S},
 \qquad
 \Gamma\longmapsto
 \bigl(\eta^D(\Gamma),\eta^\Omega(\Gamma)\bigr).
\]
\Needspace{7\baselineskip}
\subsection{Finite determination of the infinite tower}
The periodicity of the temporal word gives

\[
 D_\Gamma(9(r+12))=D_\Gamma(9r).
\]
If

\[
 T_\Gamma=\sum_{j=0}^{11}\tau_{\Gamma,j},
 \qquad r=12q+s,\quad 0\le s<12,
\]
then

\[
 A_\Gamma(12q+s)
 =A_\Gamma(s)+(12q^2+2qs)T_\Gamma^2.
\]
Therefore, the entire tower is determined by twelve values of \(D_\Gamma\), twelve values of \(A_\Gamma\), and the integer \(T_\Gamma^2\). The infinitude of heights does not require an infinitude of independent data.

The same formula solves the genealogical diamond of height 360. If

\[
 h(a,b)=90a+120b,
\]
then

\[
 (a+4,b)\sim(a,b+3)
\]
because \(3(a+4)+4b=3a+4(b+3)\). The genealogy \((a,b)\) is conserved until memory and carry are applied; only afterwards does it descend to the common height.

\Needspace{7\baselineskip}
\subsection{Absolute convergence of the refined character on the towers \(90/120\)}
We have

\[
 |\eta_h^D|\le108,\qquad
 |\eta_{30r}^{\Omega}|\le24r^2.
\]
Since \(0<q_+<q_-<1\),

\[
 |s_{30r}|\le2q_-^{30r}.
\]
It follows from this.

\[
 \sum_{h\in\mathfrak S}|\eta_h^D s_h|<\infty,
 \qquad
 \sum_{h\in\mathfrak S}|\eta_h^\Omega s_h|<\infty.
\]
The refined character is well defined without truncating the tower.

\Needspace{7\baselineskip}
\subsection{Memory, carry, and cocycle}
The bare coefficient is enriched by means of

\[
 \mathbb F_{\rm tower}(\Gamma)_h
 =\bigl(\eta_h^D,\eta_h^\Omega,\rho_h,
 \mathbf w_h,c_h,N_h^{\rm surv}\bigr),
\]
where the truncated signature satisfies

\[
 \sigma_{\Gamma|\ell_h}
 =r(\rho_h)+\mathcal D\mathbf w_h,
 \qquad
 \mathcal D=\operatorname{diag}(2,6,2,2,4),
\]
and the carry of two residues is

\[
 c(\rho_1,\rho_2)=
 \mathcal D^{-1}
 \bigl[r(\rho_1)+r(\rho_2)-r(\rho_1+\rho_2)\bigr].
\]
The locality defect of a composition,

\[
 \delta_h(\Gamma,\Lambda)=
 F_h(\Gamma\star\Lambda)-F_h(\Gamma)-F_h(\Lambda),
\]
obeys the cocycle identity

\[
 \delta_h(\Gamma,\Lambda)
 +\delta_h(\Gamma\star\Lambda,\Xi)
 =\delta_h(\Lambda,\Xi)
 +\delta_h(\Gamma,\Lambda\star\Xi).
\]
Thus it is preserved why two paths with the same visible height can carry distinct stories.

\Needspace{7\baselineskip}
\subsection{Selection of the read by coupling}
The two branches are not averaged for convenience. The physical coupling \(a\) must determine a natural transformation

\[
 \Pi_a:
 \mathcal E_{\mathfrak S}\times\mathcal E_{\mathfrak S}
 \longrightarrow\mathcal E_{\mathfrak S}
\]
that commutes with truncations and symmetries, respects the carry cocycle, preserves the electronic origin, and is compatible with \(X_{90}^{4}=Y_{120}^{3}\). If the coupling exchanges the two readings, the only continuous, multiplicative, symmetric, and normalized closure on positive scalars is the geometric mean; in exponents,

\[
 L_{\rm sym}=\frac{L_D+L_\Omega}{2}.
\]
In the absence of that symmetry, the output is the joint spectral measure of \((L_D,L_\Omega)\), not an arbitrary mean.

\Needspace{7\baselineskip}
\subsection{Orbit-primitive coinductive refinement}
The preceding finite lift admits a prolongation. Let \(\mathscr X_\Gamma\) be the graph of admissible extensions of the route, \(U_\Gamma\) its transfer operator, and \(R_\Gamma\) the orientation involution. The traces

\[
 a_n(\Gamma)=\operatorname{Tr}(R_\Gamma U_\Gamma^n)
\]
and the Möbius inversion

\[
 p_n(\Gamma)=\frac1n\sum_{d\mid n}\mu(d)a_{n/d}(\Gamma)
\]
separate primitive orbits. The difference

\[
 p_h(\Gamma)-p_h(\Gamma_e)
\]
is a further refinement of the height coefficient \(h\), not a replacement for the two finite readings.

\Needspace{7\baselineskip}
\subsection{Historical formulation via extension automaton}
The signature \(\mathbb Z^5\) does not exhaust the history of a route. For each \(\Gamma\in\mathcal R_{324}\), the following are preserved:

\begin{itemize}
\item the 108 transitions of the genealogical register;
\item the twenty blocks and their prefixes;
\item the palabras \(w_6\);
\item the nine comparisons \(K\leftrightarrow K+9\);
\item the leaf, the carry, and the winding;
\item the first returns of class, orbit, cell, and kinematic state;
\item the accumulated ambiguity and the stabilized decimal triads;
\item the profiles \(K_D\) and \(K_\Omega\).
\end{itemize}
These data allow the construction of a tower transducer without opening the residual of a mass.

\Needspace{7\baselineskip}
\subsection{Local extension automaton}
Let \(\mathscr X_\Gamma\) be the finite set of local states compatible with the route. A state preserves

\[
 x=(t,g,B,p,\varepsilon,c,w),
\]
where \(t\) is discrete time, \(g\) the nonadic phase, \(B\) the block, \(p\) the prefix, \(\varepsilon\) the sheet, \(c\) the carry and \(w\) the winding. An edge \(x\to y\) exists when \(y\) is an admissible extension of \(x\), respects the genealogical record and passes the pruning \(G_9\). The table of 36.864 carry compositions provides the local addition law; the return tables fix the memory.

Let \(U_\Gamma\) be the normalized transfer operator of this graph and \(R_\Gamma\) the orientation involution. It is required that

\[
 U_{C_M\Gamma}=J U_\Gamma J^{-1},\qquad
 R_{C_M\Gamma}=-J R_\Gamma J^{-1}.
\]
\Needspace{7\baselineskip}
\subsection{Primitive cell oriented counts and construction of the electronic signature}
For \(n\ge1\), the oriented trace

\[
 a_n(\Gamma)=\operatorname{Tr}(R_\Gamma U_\Gamma^n)
\]
counts signed closed returns. The Möbius inversion separates primitive orbits:

\[
 p_n(\Gamma)
 =\frac1n\sum_{d\mid n}\mu(d)a_{n/d}(\Gamma).
\]
The electron is taken as the affine origin of the relative corrections and is defined

\[
 \boxed{
 \eta_h(\Gamma)=p_h(\Gamma)-p_h(\Gamma_e),
 \qquad h\in\mathfrak S.}
\]
Thus the map is obtained.

\[
 \mathcal F_{\rm tower}:
 \Gamma^{\rm enr}\longmapsto
 \eta(\Gamma)=(\eta_h(\Gamma))_{h\in\mathfrak S}.
\]
There are no species-specific parameters. The same automaton, the same involution, and the same Möbius inversion act on the 324 routes.

\Needspace{7\baselineskip}
\subsection{Convergence of primitive refinement}
Let \(N_\Gamma=\dim\mathscr X_\Gamma\). If \(U_\Gamma\) is a contraction,

\[
 |a_n(\Gamma)|\le N_\Gamma.
\]
The number of divisors of \(n\) grows sublinearly, so

\[
 |p_n(\Gamma)|
 \le\frac{N_\Gamma\tau(n)}n.
\]
Since \(0<q_+<q_-<1\),

\[
 |s_h|\le q_-^h+q_+^h\le2q_-^h.
\]
By comparison,

\[
 \sum_{h\in\mathfrak S}|\eta_h(\Gamma)s_h|<\infty
\]
uniformly on any finite collection of fibers. The refined character is thus well defined.

\Needspace{7\baselineskip}
\subsection{Electronic route conservation identities and uniqueness of the central fixed point}
The construction must simultaneously satisfy:

\begin{enumerate}
\item invariance under renumbering of states;
\item covariance under leaf conjugation;
\item compatibility with direct sums of components;
\item equality upon recomputation from the genealogical register and from the certified graph;
\item conservation of the origin \(\eta(\Gamma_e)=0\);
\item absolute insensitivity to a permutation of the external mass columns;
\item determination of the tails by a bound established before comparison.
\end{enumerate}
The exact representation of a given residual proves completeness of the codomain, but does not replace these seven proofs.

\Needspace{7\baselineskip}
\section{Route clock and absolute normalization}
\Needspace{7\baselineskip}
\subsection{Integer returns are insufficient to distinguish species}
In the atlas of 324 routes, the first cell return is 27, the kinematic return within CT--108 is 54, and the extended kinematic return is 216. Being common, these integers fix the global temporal architecture, but by themselves they cannot explain mass differences between species.

\Needspace{7\baselineskip}
\subsection{Temporal holonomy of the electronic route and memory return}
The specific information resides in the holonomy of the extension operator. Let \(\lambda_j(\Gamma)\) be a nontrivial eigenvalue of \(U_\Gamma\) on the persistent subspace. We write

\[
 \lambda_j(\Gamma)=r_j(\Gamma)e^{i\theta_j(\Gamma)}.
\]
The associated internal frequency is

\[
 \omega_j(\Gamma)
 =\frac{\theta_j(\Gamma)+2\pi w_j(\Gamma)}{108t_0},
\]
where the integer \(w_j\) is fixed by the recorded winding, not by a subsequent choice of branch. For a stable mode \(r_j=1\); for a resonance \(r_j<1\), and its attenuation contributes a width.

The ratio

\[
 \rho_{\omega,j}(\Gamma)
 =\frac{\omega_j(\Gamma)}{\omega_0}
\]
is dimensionless and can modify mass ratios without altering the common decade \(-34\). The energy and mass of the mode are

\[
 E_j(\Gamma)=\hbar_{\rm ret}\omega_j(\Gamma)
 \chi_{\rm full}(\Gamma)R_{\rm act}^{\kappa_j(\Gamma)},
\]
\[
 m_j(\Gamma)=E_j(\Gamma)c_{\rm int}^{-2}.
\]
This equation brings together action, clock, character, and reader without converting any of them into another. Its numerical application requires determining \(U_\Gamma\), fixing the winding, and declaring the dimensional chart.

\Needspace{7\baselineskip}

